% ============================================================
\section{Field-Induced Horizons}
\label{sec:field-induced-horizons}
% ============================================================

The original calculus begins from a discrete horizon tower. A second and more
flexible source of horizons comes from self-adjoint fields via spectral
projection. This lets observation be defined by an internal field cut rather
than by an externally fixed scale index.

\subsection{Spectral horizon projections}
\label{subsec:field-induced-horizons}

\begin{definition}[Field-induced horizon]
\label{def:field-induced-horizon}\lean{fieldInducedHorizon}
Let $\CCField$ be a self-adjoint operator on $\cH$, and let $S \subseteq \bbR$
be a spectral window. The corresponding \emph{field-induced horizon} is the
orthogonal projection
\[
P_S := \CCSpProj{\CCField}{S} = \mathbf{1}_S(\CCField).
\]
\end{definition}

\begin{proposition}[Nested spectral cuts]
\label{prop:field-induced-nesting}\lean{fieldInducedHorizon_nesting}
If $S \subseteq T$, then
\[
P_S P_T = P_T P_S = P_S.
\]
Thus any increasing family of spectral windows determines a nested horizon
family.
\end{proposition}

\begin{proof}
Spectral projections for a self-adjoint operator commute, and the smaller
spectral window projects onto a subspace of the larger one.
\end{proof}

\begin{definition}[Field-adapted observed operator]
\label{def:field-adapted-observable}\lean{fieldAdaptedObservable}
Given an operator $A$ and a field-induced horizon $P_S$, define the observed
operator by
\[
A_{P_S} := P_S A P_S.
\]
All of the block constructions of the calculus then apply with
$P = P_S$ and $\CCComp = I - P_S$.
\end{definition}

\paragraph{Relation to the discrete tower.}
The discrete tower remains the cleanest setting for the finite proof spine.
Field-induced horizons do not replace it. They enlarge the interface of the
calculus by showing that the same observable/shadow split can arise from a
spectral selection rule.

\begin{quote}\small
\textbf{Interpretation.}
In applications, the field $\CCField$ may encode an energy, frequency,
transport, or ordering observable. The horizon is then not ``what level $h$
can be seen,'' but ``which spectral sector of $\CCField$ is retained.'' The
same defect and elimination grammar survives unchanged.
\end{quote}
